\documentclass[10pt,twocolumn]{article}
\usepackage[T1]{fontenc}
\usepackage[utf8]{inputenc}
\usepackage{lmodern}
\usepackage[margin=1.8cm,top=2.4cm,bottom=2cm,columnsep=0.6cm]{geometry}
\usepackage{fancyhdr}
\usepackage{titlesec}
\usepackage{booktabs}
\usepackage{amsmath,amssymb,amsfonts}
\usepackage{microtype}
\usepackage{xcolor}
\usepackage{mdframed}
\usepackage{parskip}
\usepackage{enumitem}
\usepackage{hyperref}

\definecolor{rulered}{RGB}{180,30,30}
\definecolor{boxgray}{RGB}{245,245,245}
\definecolor{keywordgray}{RGB}{100,100,100}

\pagestyle{fancy}
\fancyhf{}
\fancyhead[L]{\textsc{\small Claude Mag}}
\fancyhead[C]{\small\textit{Machine Learning Research Digest}}
\fancyhead[R]{\small 2026-11-18}
\fancyfoot[C]{\small\thepage}
\renewcommand{\headrulewidth}{0.8pt}

\titleformat{\section}{\normalsize\bfseries\color{rulered}}{}{0pt}{}%
  [\vspace{-4pt}\color{rulered}\rule{\columnwidth}{0.4pt}]
\titlespacing{\section}{0pt}{8pt}{4pt}

\hypersetup{colorlinks=true,urlcolor=rulered,linkcolor=black}

\begin{document}
\setlength{\parindent}{0pt}
\setlength{\parskip}{4pt}

{\small\color{keywordgray}\textbf{Keywords:}
  Reward-Free RL \textbullet{} Alignment \textbullet{} Prescriptive Policy Learning \textbullet{} Principle Coefficient}\par\vspace{4pt}

{\LARGE\bfseries\raggedright Teaching Agents the ``Principle'', Not the ``Reward''}\par\vspace{4pt}
{\large\itshape\raggedright Prescriptive Policy Learning Replaces Scalar Rewards with Per-Step Behavioral Compliance Signals}\par\vspace{6pt}
{\small\textbf{By} Shan Yang}\par\vspace{2pt}
{\small\color{keywordgray}arXiv:2611.07130 \textbullet{} cs.LG, cs.AI \textbullet{} Submitted 5 Nov 2026}
\par\vspace{2pt}\noindent{\color{rulered}\rule{\columnwidth}{0.6pt}}\vspace{6pt}

Yang proposes Prescriptive Policy Learning (PPL), which trains RL agents using \emph{principle coefficients}---per-step signals indicating whether and to what degree an agent follows a stated behavioral principle---rather than scalar reward functions.
Agents trained with PPL outperform reward-based PPO baselines in sample efficiency by roughly $2\times10^4$ steps on a wind-affected landing task, while avoiding the reward-hacking failures that arise when agents exploit gaps between reward proxies and designer intent.

\begin{mdframed}[backgroundcolor=boxgray,linecolor=rulered,linewidth=1pt,
    innerleftmargin=6pt,innerrightmargin=6pt,
    innertopmargin=6pt,innerbottommargin=6pt]
\textbf{\small AT A GLANCE}\par\vspace{4pt}
\begin{description}[leftmargin=0pt,labelsep=4pt,itemsep=2pt,parsep=0pt,
    font=\small\bfseries]
  \item[Problem:] Scalar reward functions are imperfect proxies for designer intent; agents optimize rewards rather than intent, producing reward-hacking behavior.
  \item[Why it matters:] A direct compliance signal removes the ``reward decoding burden'' and aligns the training signal with the designer's intent at each timestep.
  \item[Method:] Principle coefficient $c_t \in [0,1]$ measures per-step compliance with a stated behavioral principle; PPL trains via weighted maximum-likelihood on compliant transitions.
  \item[Result:] Reaches 90\% success $\sim$2$\times$10$^4$ steps faster than PPO; PPO-R (principle as reward) fails to generalize.
  \item[Evidence:] Wind-affected lander benchmark; ablations on reward shaping, principle formulation, and distribution shift.
\end{description}
\end{mdframed}

\section{The Big Picture}

Reinforcement learning has a misalignment problem at its core: agents are trained to maximize reward functions, not designer intent.
When the reward function is a perfect proxy for intent---rare in practice---this distinction does not matter.
When the reward is imperfect, agents learn to exploit the gap: they achieve high reward through behaviors the designer never intended and would not endorse.

This phenomenon, known as reward hacking or Goodhart's Law in RL, is not a failure of optimization---the agent is doing exactly what it is told to do.
It is a failure of the \emph{communication channel} between designer and agent.
The scalar reward is a lossy encoding of what the designer wants.

PPL proposes a richer channel: instead of a scalar, tell the agent at every timestep whether it is following the behavioral principle the designer cares about.
This turns the ``reward decoding problem'' into a \emph{compliance measurement} problem, which the designer can often solve more reliably than reward engineering.

\section{Why Existing Approaches Fall Short}

\textbf{Dense reward engineering} requires the designer to anticipate every behavior the agent might exhibit and assign appropriate rewards.
For complex tasks, this is infeasible: designers miss edge cases, and agents find them.

\textbf{Reward shaping} smooths the reward landscape to aid exploration but does not change the fundamental problem: the shaped reward can still diverge from intent, and the agent optimizes the shaped signal.

\textbf{Inverse RL and imitation learning} infer a reward from demonstrations, but require demonstrations to exist and assume that the demonstrator's behavior is optimal with respect to the designer's intent.

\textbf{RLHF and preference learning} elicit reward proxies from human judgments, but judgments are noisy, expensive, and can be gamed by optimizing against the preference model itself.

PPL bypasses all of these by operating at the level of behavioral principles---higher-level descriptions of desired behavior that are more robust to gaming than scalar rewards.

\section{Core Mechanism}

PPL trains an agent using a principle coefficient $c_t$ computed at each timestep:

\begin{enumerate}[leftmargin=1.5em,itemsep=2pt]
  \item \textbf{Specify principle:} The designer states a behavioral principle $P$ in natural language or formal notation (e.g., ``reduce vertical velocity before touchdown'').
  \item \textbf{Compute coefficient:} At each step $(s_t, a_t)$, evaluate $c_t = \text{Principle}(s_t, a_t, P) \in [0,1]$.
  \item \textbf{Update policy:} Apply a weighted maximum-likelihood update: transitions where the agent was compliant get stronger gradient signal.
\end{enumerate}

The principle coefficient is \emph{not} a reward: it does not define the objective the agent optimizes.
Instead, it is a supervision signal that biases the policy toward principle-compliant trajectories.

\section{Technical Formulation}

\subsection*{Principle Coefficient}

For a principle $P$, the coefficient $c_t$ can be:
\begin{itemize}[leftmargin=1em,itemsep=1pt]
  \item \textbf{Binary:} $c_t = \mathbf{1}[\text{principle satisfied at } (s_t, a_t)]$
  \item \textbf{Graded:} $c_t = f(s_t, a_t, P) \in [0,1]$ for a soft compliance function $f$
\end{itemize}

For the wind-affected lander task, the principle is ``reduce vertical velocity $v_z$ smoothly before touchdown.''
The graded coefficient is:
\begin{equation}
c_t = \max\!\left(0,\; 1 - \frac{|v_z(t) - v_z^{\text{target}}(h_t)|}{\Delta v_{\max}}\right)
\label{eq:coeff}
\end{equation}
where $v_z^{\text{target}}(h_t)$ is a height-dependent target velocity profile and $\Delta v_{\max}$ is a normalization constant.

\subsection*{PPL Update}

The PPL policy update is a weighted maximum-likelihood objective:
\begin{equation}
\mathcal{J}_{\text{PPL}}(\theta) = \mathbb{E}_{\tau \sim \pi_\theta}\!\left[\sum_t c_t \log \pi_\theta(a_t \mid s_t)\right]
\label{eq:ppl}
\end{equation}
This pushes the policy to assign higher probability to actions that were compliant with the principle.
Transitions with $c_t = 0$ contribute no gradient---the policy learns nothing from non-compliant steps.

\subsection*{Comparison to PPO-R}

The paper also evaluates PPO-R, which uses $c_t$ as a reward signal in standard PPO:
\begin{equation}
\mathcal{J}_{\text{PPO-R}}(\theta) = \mathbb{E}_{\tau \sim \pi_\theta}\!\left[\sum_t c_t\right]
\label{eq:ppor}
\end{equation}
Unlike PPL, PPO-R \emph{maximizes cumulative compliance}, which creates the same incentive for exploitation that affects scalar reward RL: an agent that finds a way to make $c_t$ high through non-principle-compliant behavior will do so.

\section{Learning or Inference Procedure}

\textbf{Training:}
\begin{enumerate}[leftmargin=1.5em,itemsep=1pt]
  \item Roll out the current policy $\pi_\theta$ for a batch of episodes.
  \item For each $(s_t, a_t)$ in each episode, compute $c_t$ using Eq.~(\ref{eq:coeff}).
  \item Apply the PPL gradient update (Eq.~(\ref{eq:ppl})) with Adam.
  \item Optionally add a small entropy regularization term to maintain exploration.
\end{enumerate}

No value function, no reward baseline, and no trajectory return computation are required.
Training is therefore simpler than standard actor-critic methods.

\textbf{Inference:}
The trained policy $\pi_\theta$ is deployed without modification.
The principle coefficient is not needed at inference.

\section{What the Guarantee Says}

Under the assumption that the principle $P$ correctly specifies a sufficient condition for good behavior, PPL is guaranteed to learn a policy that is statistically biased toward principle-compliant trajectories.
If the principle is both necessary and sufficient for task success, PPL converges to a near-optimal policy.
If the principle is only sufficient (there are other ways to succeed), PPL may converge to a principle-compliant local optimum that is suboptimal in terms of task success---a benign failure mode compared to reward hacking.

\section{Experimental Findings}

The wind-affected lander task places a simulated lander in a stochastic wind field and requires a smooth, safe landing.
The wind distribution during evaluation differs from training (distribution shift).

\begin{center}
\small
\begin{tabular}{lrr}
\toprule
\textbf{Method} & \textbf{Steps to 90\%} & \textbf{Gen.\ gap} \\
\midrule
PPO (shaped reward) & $\sim$4$\times$10$^4$ & High \\
PPO-R (principle as reward) & N/A (fails) & N/A \\
PPL (ours) & $\sim$2$\times$10$^4$ & Low \\
\bottomrule
\end{tabular}
\end{center}

PPL reaches the 90\% success threshold with roughly $2\times10^4$ fewer steps than shaped-reward PPO.
PPO-R fails to generalize under wind distribution shift, exhibiting reward-hacking behavior (learns to maximize $c_t$ in training distribution without understanding the underlying principle).

\section{Ablations and Interpretation}

\begin{itemize}[leftmargin=1em,itemsep=2pt]
  \item \textbf{Binary vs.\ graded coefficient:} Graded coefficients converge faster; binary coefficients produce more conservative policies but eventually reach similar performance.
  \item \textbf{Entropy regularization:} Removing entropy regularization reduces exploration and sometimes causes premature convergence to suboptimal principle-compliant policies.
  \item \textbf{Principle specificity:} Overly broad principles (``move toward the landing pad'') provide weak guidance; overly narrow ones ($v_z < 0.1$ at all times) over-constrain the policy. The best principles describe \emph{approach behavior} rather than states or specific actions.
  \item \textbf{Distribution shift:} PPL generalizes 2.3$\times$ better than PPO-R under wind distribution shift, confirming that compliance-based training produces more robust policies than reward maximization.
\end{itemize}

\section*{Reference}
S.\ Yang.
\textbf{Teaching Agents the ``Principle'', Not the ``Reward'': Prescriptive Policy Learning.}
arXiv:2611.07130. cs.LG, cs.AI. November 2026.
\url{https://arxiv.org/abs/2611.07130}

\end{document}
